\section{Introdução}
\subsection{Posicionamento e Declaração do Problema}
% [Forneça uma declaração resumindo o problema que está sendo resolvido por este projeto.]
O problema de \textbf{[descreva o problema]} afeta \textbf{[as partes interessadas afetadas pelo problema]}. O impacto disso é \textbf{[qual é o impacto do problema?]}, e uma solução bem-sucedida seria \textbf{[liste alguns benefícios principais de uma solução bem-sucedida]}.

\subsection{Declaração de Posição do Produto}
% [Forneça uma declaração geral resumindo a posição única que o produto pretende ocupar no mercado.]
\begin{itemize}
    \item \textbf{Para:} [cliente alvo]
    \item \textbf{Que:} [declaração da necessidade ou oportunidade]
    \item \textbf{O (nome do produto):} é uma [categoria do produto]
    \item \textbf{Que:} [declaração do benefício principal; ou seja, a razão convincente para usar/comprar]
    \item \textbf{Diferente de:} [principal alternativa competitiva]
    \item \textbf{Nosso produto:} [declaração da diferenciação principal]
\end{itemize}

\section{Descrições das Partes Interessadas (Stakeholders)}
% [Resuma os tipos de stakeholders e suas responsabilidades]
\begin{table}[H]
    \centering
    \begin{tabular}{|p{3cm}|p{5cm}|p{6cm}|}
        \hline
        \textbf{Nome} & \textbf{Descrição} & \textbf{Responsabilidades} \\ \hline
        [Stakeholder 1] & [Descreva brevemente o stakeholder] & [Resuma as principais responsabilidades. Ex: Garante que o sistema atenda ao mercado, aprova financiamento, etc.] \\ \hline
        [Stakeholder 2] & [Descreva] & [Descreva] \\ \hline
    \end{tabular}
    \caption{Resumo das Partes Interessadas}
\end{table}

\section{Ambiente do Usuário}
% [Detalhe o ambiente de trabalho do usuário alvo. Sugestões: Quantas pessoas? Ciclo da tarefa? Restrições de ambiente (mobile, etc.)? Plataformas em uso? Integração com outros apps?]
\begin{itemize}
    \item \textbf{Número de pessoas envolvidas nas tarefas:} [Detalhar]
    \item \textbf{Ciclo e tempo das atividades:} [Detalhar]
    \item \textbf{Restrições do ambiente:} [Ex: Mobile, web, offline]
    \item \textbf{Plataformas do sistema em uso:} [Ex: Windows, Linux, iOS]
    \item \textbf{Integrações necessárias:} [Ex: Sistema legado, ERP]
\end{itemize}

\section{Visão Geral do Produto}
\subsection{Necessidades e Recursos (Features)}
% [Concentre-se nos recursos necessários e por que devem ser implementados. Evite design aqui. Capture prioridade.]
\begin{table}[H]
    \centering
    \begin{tabular}{|p{4cm}|p{2cm}|p{6cm}|p{2cm}|}
        \hline
        \textbf{Necessidade} & \textbf{Prioridade} & \textbf{Recursos (Features)} & \textbf{Lançamento Planejado} \\ \hline
        [Descreva a necessidade] & [Alta/Média] & [Descreva a funcionalidade que atende à necessidade] & [Versão/Data] \\ \hline
    \end{tabular}
    \caption{Necessidades e Funcionalidades}
\end{table}

\subsection{Outros Requisitos do Produto}
% [Em alto nível, liste padrões aplicáveis, hardware, requisitos de desempenho, ambientais, limites de projeto (design constraints), documentação, etc.]
\begin{table}[H]
    \centering
    \begin{tabular}{|p{6cm}|p{2cm}|p{6cm}|}
        \hline
        \textbf{Requisito} & \textbf{Prioridade} & \textbf{Lançamento Planejado} \\ \hline
        [Ex: Requisitos de usabilidade e performance] & [Alta] & [Versão 1.0] \\ \hline
        [Ex: Padrões e leis aplicáveis] & [Alta] & [Versão 1.0] \\ \hline
        [Ex: Requisitos de documentação (manuais)] & [Média] & [Versão 1.0] \\ \hline
    \end{tabular}
    \caption{Outros Requisitos do Produto}
\end{table}
